% ============================================================================
% Appendix: Implementation Details
% ============================================================================

\section{Implementation Details}
\label{app:impl}

This appendix provides implementation details for DUET beyond the main method description.
We focus on algorithmic choices needed for reproducibility, and omit code-level configuration names.

\subsection{Density-ratio discriminator}
\label{app:disc}

DUET estimates the teacher--student density ratio with a lightweight discriminator.
The discriminator is a three-layer MLP with hidden width 64 and $\tanh$ activations.
Its input consists of sequence-level features computed from each rollout, including summary statistics of token log-probabilities, low-probability tail ratios, KL-to-reference statistics, and response length.
When high-dimensional hidden features are used, we first apply a LayerNorm projection with dropout before concatenating them with scalar features.

The discriminator is trained with binary cross-entropy to distinguish on-policy samples from teacher-replay samples.
Training batches are class-balanced between the two sources, and we use label smoothing with coefficient $0.1$.
The optimizer is Adam with learning rate $3\times 10^{-4}$ and weight decay $0.005$.
Discriminator examples are drawn from a rolling buffer of recent on-policy and teacher samples.
To emphasize recent policy behavior while retaining stability, older samples are down-weighted exponentially by age and the weights are normalized to preserve the loss scale.

We delay applying the density-ratio weights until the discriminator has accumulated a sufficient buffer of examples.
Before this warmup completes, teacher samples are trained without DR3 weighting and behavior cloning remains at its cold-start strength.
This avoids applying unreliable density-ratio estimates during the first few actor updates.

For inference, we use a smoothed copy of the discriminator parameters to compute $\hat w$.
This reduces short-term feedback between discriminator updates and the actor loss that consumes the resulting weights.
We also apply a temperature to discriminator logits before converting them into density ratios, which prevents over-confident early predictions from producing unstable replay weights.

\subsection{Replay-weight clipping and stabilization}
\label{app:ratio_clip}

The raw density-ratio estimate
\[
\hat w(s,a)=\frac{D_\phi(s,a)}{1-D_\phi(s,a)}
\]
can be noisy early in training.
We therefore clip the effective teacher replay multiplier used in the actor loss.
The clipping threshold is adapted to keep the effective sample size of teacher weights within a stable range.
Concretely, we maintain a dual variable that increases when the effective sample size falls below a target fraction of the teacher batch and decreases otherwise.
The resulting upper bound prevents a small number of teacher samples from dominating the policy-gradient update.

In distributed training, discriminator features and labels are synchronized across workers before discriminator updates.
The discriminator parameters and scalar statistics used by the actor loss are also synchronized across workers.
This ensures that all ranks apply consistent replay weights and behavior-cloning coefficients.

\subsection{Adaptive behavior-cloning coefficient}
\label{app:bc_schedule}

The behavior-cloning coefficient $\mu(t)$ is driven by discriminator accuracy rather than by a fixed step-indexed schedule.
Let $\bar d_t$ be an exponential moving average of discriminator accuracy, and let $d_{\rm floor}$ denote the accuracy threshold below which the discriminator is treated as uninformative.
We map discriminator accuracy to the BC coefficient with
\begin{equation}
g_t
=
\mathrm{clip}
\left(
\frac{1-\bar d_t}{1-d_{\rm floor}},
0,
1
\right),
\qquad
\mu(t)
=
\mu_{\min}
+
(\mu_{\max}-\mu_{\min})g_t .
\label{eq:appendix_bc_schedule}
\end{equation}
Thus, BC remains strong while the discriminator is near chance and decays as the discriminator becomes informative.
During discriminator warmup, we treat accuracy as chance-level so that $\mu(t)$ stays at $\mu_{\max}$.

We use token-level cross-entropy on teacher trajectories as the BC objective.
No additional token-level weighting is applied inside the BC loss, so $\mu(t)$ is the only coefficient controlling the strength of behavior cloning.
The hyperparameters used in the 1.5B experiments are shown in Table~\ref{tab:bc_hparams}.

\begin{table}[h]
\centering
\caption{
\textbf{Adaptive BC hyperparameters.}
$\mu_{\max}$ and $\mu_{\min}$ are the maximum and minimum BC coefficients.
$d_{\rm floor}$ is the discriminator-accuracy threshold used in Eq.~\eqref{eq:appendix_bc_schedule}.
}
\label{tab:bc_hparams}
\begin{tabular}{lcccc}
\toprule
Setting & $\mu_{\max}$ & $\mu_{\min}$ & $d_{\rm floor}$ & EMA decay \\
\midrule
1.5B-ALFWorld & 0.30 & 0.05 & 0.40 & 0.50 \\
1.5B-WebShop  & 0.30 & 0.10 & 0.60 & 0.20 \\
\bottomrule
\end{tabular}
\end{table}

\subsection{State Channel implementation}
\label{app:state_channel_impl}

The State Channel builds a progress potential $\Phi$ from successful teacher trajectories.
For each teacher trajectory, observations are assigned normalized progress values according to their relative position in the trajectory.
These values are aggregated across the teacher cache to form a map from environment states to progress scores in $[0,1]$.
States that cannot be matched to the teacher cache receive progress value $0$.

State matching is environment dependent.
For ALFWorld, observations are represented with a compact state key derived from the text observation.
For WebShop, the state key additionally uses product and attribute information, since progress depends on matching user-specified constraints such as product type, brand, price, and other attributes.
This makes the progress map more robust to surface-form variation in WebShop observations.

For an on-policy trajectory $\tau=(s_0,a_0,\ldots,s_T)$, the progress score is the mean potential over visited states:
\begin{equation}
P(\tau)
=
\frac{1}{T+1}\sum_{t=0}^{T}\Phi(s_t).
\label{eq:appendix_progress_score}
\end{equation}
We use $\lambda=0.20$ for the shaping coefficient and add the bonus $\lambda P(\tau)$ to the trajectory reward.
The bonus is applied only to on-policy rollouts.
Teacher trajectories keep their original success reward, since adding progress bonuses to already-successful teacher samples would increase their advantages and partially counteract DR3.

The State Channel bonus is added to the on-policy reward used by the policy-gradient update, but it is not used to shift the teacher baseline.
This keeps state-level shaping separate from the Action Channel correction.

\subsection{Asymmetric baseline separation}
\label{app:baseline_sep_impl}

DUET uses the asymmetric baseline rule described in Section~\ref{sec:action_channel}.
For on-policy samples, the reward baseline is the mean reward of the on-policy subgroup:
\begin{equation}
\hat A^{o,(i)}
=
\frac{R^{o,(i)}-\mu^o}{\sigma^o}.
\label{eq:appendix_on_policy_adv}
\end{equation}
For teacher samples, the baseline is the pooled group mean:
\begin{equation}
\hat A^{\beta,(i)}
=
\frac{R^{\beta,(i)}-\mu_{\rm pool}}{\sigma^o}.
\label{eq:appendix_teacher_adv}
\end{equation}
Both subgroups are scaled by the on-policy standard deviation $\sigma^o$.
This choice avoids the zero-variance problem caused by success-filtered teacher batches, where all teacher rewards are often identical.
It also ensures that teacher rewards do not shift the on-policy baseline, while teacher samples still receive positive advantage when the student is weak.

\subsection{Training system}
\label{app:software}

Training uses a distributed PPO-style RL loop with fully sharded data-parallel policy training.
Rollout generation is served asynchronously, and environment interaction is parallelized across separate actor pools.
The student policy, rollout workers, discriminator, and reward computation are synchronized once per actor update.
All methods in the main experiments use the same training infrastructure, teacher cache, optimization budget, and evaluation protocol.